% problem_id: S001-lemma-compare-higher
% source_id: S001 (Stacks Project)
% source_file: cotangent.tex
% statement_locator: cotangent.tex lines 2452--2458
% proof_locator: cotangent.tex lines 2460--2605
% env: lemma
% id_note: label 在本来源内唯一
% pairing: tight (verified)
% status: complete (statement + author proof verbatim + reason + locators)
%
\begin{lemma}
\label{lemma-compare-higher}
In the situation above denote $L$ the complex
(\ref{equation-lichtenbaum-schlessinger}).
There is a canonical map $L_{B/A} \to L$ in $D(B)$ which
induces an isomorphism $\tau_{\geq -2}L_{B/A} \to L$ in $D(B)$.
\end{lemma}

\begin{proof}
Let $P_\bullet \to B$ be a resolution of $B$ over $A$
(Remark \ref{remark-resolution}). We will identify $L_{B/A}$ with
$\Omega_{P_\bullet/A} \otimes B$. To construct the map we
make some choices.

\medskip\noindent
Choose an $A$-algebra map $\psi : P_0 \to P$ compatible with the
given maps $P_0 \to B$ and $P \to B$.

\medskip\noindent
Write $P_1 = A[S]$ for some set $S$. For $s \in S$ we may write
$$
\psi(d_0(s) - d_1(s)) = \sum p_{s, t} f_t
$$
for some $p_{s, t} \in P$. Think of $F = \bigoplus_{t \in T} P$
as a $(P_1, P_1)$-bimodule via the maps $(\psi \circ d_0, \psi \circ d_1)$.
By Lemma \ref{lemma-polynomial-ring-unique} we obtain a unique
$A$-biderivation $\lambda : P_1 \to F$ mapping $s$ to the vector with
coordinates $p_{s, t}$. By construction the composition
$$
P_1 \longrightarrow F \longrightarrow P
$$
sends $f \in P_1$ to $\psi(d_0(f) - d_1(f))$ because the map
$f \mapsto \psi(d_0(f) - d_1(f))$ is an $A$-biderivation agreeing with
the composition on generators.

\medskip\noindent
For $g \in P_2$ we claim that $\lambda(d_0(g) - d_1(g) + d_2(g))$
is an element of $Rel$. Namely, by the last remark of the previous
paragraph the image of $\lambda(d_0(g) - d_1(g) + d_2(g))$ in $P$ is
$$
\psi((d_0 - d_1)(d_0(g) - d_1(g) + d_2(g)))
$$
which is zero by Simplicial, Section \ref{simplicial-section-complexes}).

\medskip\noindent
The choice of $\psi$ determines a map
$$
\text{d}\psi \otimes 1 :
\Omega_{P_0/A} \otimes B
\longrightarrow
\Omega_{P/A} \otimes B
$$
Composing $\lambda$ with the map $F \to F \otimes B$ gives a
usual $A$-derivation as the two $P_1$-module structures on
$F \otimes B$ agree. Thus $\lambda$ determines a map
$$
\overline{\lambda} :
\Omega_{P_1/A} \otimes B
\longrightarrow
F \otimes B
$$
Finally, We obtain a $B$-linear map
$$
q :
\Omega_{P_2/A} \otimes B
\longrightarrow
Rel/TrivRel
$$
by mapping $\text{d}g$ to the class of $\lambda(d_0(g) - d_1(g) + d_2(g))$
in the quotient.

\medskip\noindent
The diagram
$$
\xymatrix{
\Omega_{P_3/A} \otimes B \ar[r] \ar[d] &
\Omega_{P_2/A} \otimes B \ar[r] \ar[d]_q &
\Omega_{P_1/A} \otimes B \ar[r] \ar[d]_{\overline{\lambda}} &
\Omega_{P_0/A} \otimes B \ar[d]_{\text{d}\psi \otimes 1} \\
0 \ar[r] &
Rel/TrivRel \ar[r] &
F \otimes B \ar[r] &
\Omega_{P/A} \otimes B
}
$$
commutes (calculation omitted) and we obtain the map of the lemma.
By Remark \ref{remark-explicit-comparison-map} and
Lemma \ref{lemma-relation-with-naive-cotangent-complex} we see that this map
induces isomorphisms $H_1(L_{B/A}) \to H_1(L)$ and $H_0(L_{B/A}) \to H_0(L)$.

\medskip\noindent
It remains to see that our map $L_{B/A} \to L$ induces an isomorphism
$H_2(L_{B/A}) \to H_2(L)$. Choose a resolution of $B$ over $A$ with
$P_0 = P = A[u_i]$ and then $P_1$ and $P_2$ as in
Example \ref{example-resolution-length-two}.
In Remark \ref{remark-elucidate-degree-two} we have constructed an exact
sequence
$$
\wedge^2_B(J_0/J_0^2) \to \text{Tor}_2^{P_0}(B, B) \to H^{-2}(L_{B/A}) \to 0
$$
where $P_0 = P$ and $J_0 = \Ker(P \to B) = I$.
Calculating the Tor group using the short exact sequences
$0 \to I \to P \to B \to 0$ and $0 \to Rel \to F \to I \to 0$
we find that
$\text{Tor}_2^P(B, B) = \Ker(Rel \otimes B \to F \otimes B)$.
The image of the map $\wedge^2_B(I/I^2) \to \text{Tor}_2^P(B, B)$
under this identification is exactly the image of $TrivRel \otimes B$.
Thus we see that $H_2(L_{B/A}) \cong H_2(L)$.

\medskip\noindent
Finally, we have to check that our map $L_{B/A} \to L$ actually induces
this isomorphism. We will use the notation and results discussed in
Example \ref{example-resolution-length-two} and
Remarks \ref{remark-elucidate-degree-two} and \ref{remark-surjection}
without further mention. Pick an element $\xi$ of
$\text{Tor}_2^{P_0}(B, B) = \Ker(I \otimes_P I \to I^2)$.
Write $\xi = \sum h_{t', t}f_{t'} \otimes f_t$ for some
$h_{t', t} \in P$. Tracing through the exact sequences above we
find that $\xi$ corresponds to the image in $Rel \otimes B$
of the element $r \in Rel \subset F = \bigoplus_{t \in T} P$ with
$t$th coordinate $r_t = \sum_{t' \in T} h_{t', t}f_{t'}$.
On the other hand, $\xi$ corresponds to the element of
$H_2(L_{B/A}) = H_2(\Omega)$ which is the image
via $\text{d} : H_2(\mathcal{J}/\mathcal{J}^2) \to H_2(\Omega)$
of the boundary of $\xi$ under the $2$-extension
$$
0 \to
\text{Tor}_2^\mathcal{O}(\underline{B}, \underline{B})
\to 
\mathcal{J} \otimes_\mathcal{O} \mathcal{J} \to \mathcal{J}
\to
\mathcal{J}/\mathcal{J}^2 \to 0
$$
We compute the successive transgressions of our element. First we have
$$
\xi = (d_0 - d_1)(- \sum s_0(h_{t', t} f_{t'}) \otimes x_t)
$$
and next we have
$$
\sum s_0(h_{t', t} f_{t'}) x_t = d_0(v_r) - d_1(v_r) + d_2(v_r)
$$
by our choice of the variables $v$ in
Example \ref{example-resolution-length-two}.
We may choose our map $\lambda$ above such that
$\lambda(u_i) = 0$ and $\lambda(x_t) = - e_t$ where $e_t \in F$
denotes the basis vector corresponding to $t \in T$.
Hence the construction of our map $q$ above sends $\text{d}v_r$ to
$$
\lambda(\sum s_0(h_{t', t} f_{t'}) x_t) =
\sum\nolimits_t \left(\sum\nolimits_{t'} h_{t', t}f_{t'}\right) e_t
$$
matching the image of $\xi$ in $Rel \otimes B$ (the two minus signs
we found above cancel out). This agreement finishes the proof.
\end{proof}
